% Compile with latexmk (uses ~/.latexmkrc: LuaLaTeX, output in build/)
% Week 10: continuous random variables, density, distribution function,
% expected value and variance, the uniform, exponential and normal
% distributions; the distribution of a function of a random variable, joint
% distributions (discrete and continuous), independent random variables.
% Sessions 1 and 2 are held on the SAME DAY, back to back (2 x 90 min,
% short break in between). The lecture of this week is given by someone
% else: the slides assume nothing from it and teach every concept.
% Reference: Ross, A First Course in Probability, 8th ed., Ch. 5--6
% (notation and the problems paraphrased).
% The exercises are in exercises/ (one file each), placed by \exerciseframe.

\documentclass[aspectratio=169,11pt,usepdftitle=false]{beamer}
\usepackage{fontspec}
\setsansfont{Latin Modern Sans}
\usepackage{amsmath,amssymb,mathtools}
\usepackage{tikz}
\usetikzlibrary{calc,arrows.meta}
\usepackage{booktabs}
\usepackage{array}
\usepackage{cancel}

\usetheme{Madrid}

\definecolor{navy}{RGB}{24,48,75}
\definecolor{teal}{RGB}{0,120,125}
\definecolor{gold}{RGB}{190,130,20}

\setbeamercolor{structure}{fg=navy}
\setbeamercolor{alerted text}{fg=teal}
\setbeamercolor{block title}{bg=navy,fg=white}
\setbeamercolor{block body}{bg=navy!5}
\setbeamercolor{block title example}{bg=teal,fg=white}
\setbeamercolor{block body example}{bg=teal!6}
\setbeamercolor{block title alerted}{bg=gold,fg=white}
\setbeamercolor{block body alerted}{bg=gold!8}

\setbeamertemplate{navigation symbols}{}
\setbeamertemplate{footline}{%
  \usebeamerfont{footline}\hspace{1em}\insertshortauthor\hfill%
  \insertshorttitle\quad|\quad%
  \insertframenumber/\inserttotalframenumber%
  \hspace{1em}\vspace{2mm}%
}
\setbeamertemplate{itemize items}[circle]

% Contact (beamer loads hyperref itself)
\newcommand{\myemail}{tiborsandor2001@gmail.com}

\newcommand{\N}{\mathbb{N}}
\newcommand{\Om}{\Omega}
% Expectation, variance and standard deviation, as in Ross: E[X],
% Var(X), SD(X)
\newcommand{\E}{\operatorname{E}}
\newcommand{\Var}{\operatorname{Var}}
\newcommand{\SD}{\operatorname{SD}}
\newcommand{\sol}{\alert{Solution.}\ }

% Pictures of continuous distributions (teal areas, navy curves, gold marks).
% The same definitions are in beamer.tex and preamble.tex.
% A density f(x) = x e^{-x} (x >= 0) and the area P(a <= X <= b) under it,
% a = 1.5, b = 3.2 (area 0.387).
\newcommand{\figdensity}{%
  \begin{tikzpicture}[x=.85cm,y=5cm]
    \fill[teal!25] (1.5,0) -- plot[domain=1.5:3.2,samples=40]
      (\x,{\x*exp(-\x)}) -- (3.2,0) -- cycle;
    \draw[->] (-.2,0)--(6.4,0) node[right,font=\scriptsize] {$x$};
    \draw[->] (0,-.02)--(0,.47);
    \draw[very thick,navy] (-.2,0)--(0,0) plot[domain=0:6.1,samples=80]
      (\x,{\x*exp(-\x)});
    \node[navy,font=\scriptsize] at (1.7,.42) {$f(x)$};
    \node[below,font=\scriptsize] at (1.5,0) {$a$};
    \node[below,font=\scriptsize] at (3.2,0) {$b$};
    \draw[teal,->] (4.1,.24)--(2.5,.13);
    \node[right,font=\scriptsize,align=left] at (3.9,.29)
      {area $=$\\$P(a\le X\le b)$};
  \end{tikzpicture}}
% The distribution function F(x) = 1 - (1+x) e^{-x} of the same density;
% the arrow marks F(b) - F(a), the same area.
\newcommand{\figcdf}{%
  \begin{tikzpicture}[x=.8cm,y=1.9cm]
    \draw[->] (-.2,0)--(6.4,0) node[right,font=\scriptsize] {$x$};
    \draw[->] (0,-.05)--(0,1.2);
    \draw[gray!50,dashed] (0,1)--(6.1,1);
    \node[left,font=\scriptsize] at (0,1) {$1$};
    \draw[very thick,navy] (-.2,0)--(0,0) plot[domain=0:6.1,samples=80]
      (\x,{1-(1+\x)*exp(-\x)});
    \node[navy,font=\scriptsize] at (1.0,.62) {$F(x)$};
    \draw[gold,densely dotted] (1.5,0)--(1.5,0.442)--(6.1,0.442)
      (3.2,0)--(3.2,0.829)--(6.1,0.829);
    \node[below,font=\scriptsize] at (1.5,0) {$a$};
    \node[below,font=\scriptsize] at (3.2,0) {$b$};
    \draw[gold,very thick,<->] (5.6,0.442)--(5.6,0.829);
    \node[right,font=\scriptsize] at (5.6,.635) {$F(b)-F(a)$};
  \end{tikzpicture}}
% Uniform density on [alpha, beta]; the shaded interval [c, d] has
% probability (d - c)/(beta - alpha); the gold mark is E[X].
\newcommand{\figuniform}{%
  \begin{tikzpicture}[x=.9cm,y=1.6cm]
    \fill[teal!25] (2.1,0) rectangle (3.1,.8);
    \draw[->] (-.2,0)--(5.4,0) node[right,font=\scriptsize] {$x$};
    \draw[->] (0,-.05)--(0,1.25);
    \draw[very thick,navy] (-.2,0)--(1,0) (1,.8)--(4.6,.8) (4.6,0)--(5.2,0);
    \draw[navy,densely dotted] (1,0)--(1,.8) (4.6,0)--(4.6,.8);
    \draw[gray!50,dashed] (0,.8)--(1,.8);
    \node[left,font=\scriptsize] at (0,.8) {$\frac{1}{\beta-\alpha}$};
    \node[below,font=\scriptsize] at (1,0) {$\alpha$};
    \node[below,font=\scriptsize] at (4.6,0) {$\beta$};
    \node[below,font=\scriptsize] at (2.1,0) {$c$};
    \node[below,font=\scriptsize] at (3.1,0) {$d$};
    \node[font=\tiny] at (2.6,.4) {$\frac{d-c}{\beta-\alpha}$};
    \draw[gold,very thick,dashed] (2.8,.8)--(2.8,1.12);
    \node[above,font=\tiny] at (2.8,1.1) {$\E[X]=\frac{\alpha+\beta}{2}$};
  \end{tikzpicture}}
% Exponential density with rate lambda (drawn with lambda = 1): the shaded
% tail beyond x has area e^{-lambda x}; the gold mark is E[X] = 1/lambda.
\newcommand{\figexponential}{%
  \begin{tikzpicture}[x=.9cm,y=1.5cm]
    \fill[teal!25] (1.6,0) -- plot[domain=1.6:5,samples=40]
      (\x,{exp(-\x)}) -- (5,0) -- cycle;
    \draw[->] (-.2,0)--(5.4,0) node[right,font=\scriptsize] {$t$};
    \draw[->] (0,-.05)--(0,1.25);
    \draw[very thick,navy] (-.2,0)--(0,0);
    \draw[very thick,navy] plot[domain=0:5,samples=60] (\x,{exp(-\x)});
    \node[left,font=\scriptsize] at (0,1) {$\lambda$};
    \node[navy,font=\scriptsize] at (1.05,.75) {$\lambda e^{-\lambda t}$};
    \draw[gold,very thick,dashed] (1,0)--(1,.37);
    \node[below,font=\scriptsize] at (1,0) {$\frac1\lambda$};
    \node[below,font=\scriptsize] at (1.6,0) {$x$};
    \draw[teal,->] (3.3,.55)--(2.3,.08);
    \node[above,font=\scriptsize] at (3.6,.52) {$P(X>x)=e^{-\lambda x}$};
  \end{tikzpicture}}
% Normal density with mean mu and standard deviation sigma: the bands
% mu +- sigma, mu +- 2 sigma, mu +- 3 sigma and their probabilities.
\newcommand{\fignormal}{%
  \begin{tikzpicture}[x=.8cm,y=5.2cm]
    \fill[teal!10] (-3,0) -- plot[domain=-3:3,samples=80]
      (\x,{exp(-\x*\x/2)/sqrt(2*pi)}) -- (3,0) -- cycle;
    \fill[teal!22] (-2,0) -- plot[domain=-2:2,samples=60]
      (\x,{exp(-\x*\x/2)/sqrt(2*pi)}) -- (2,0) -- cycle;
    \fill[teal!40] (-1,0) -- plot[domain=-1:1,samples=40]
      (\x,{exp(-\x*\x/2)/sqrt(2*pi)}) -- (1,0) -- cycle;
    \draw[->] (-3.7,0)--(3.9,0) node[right,font=\scriptsize] {$x$};
    \draw[very thick,navy] plot[domain=-3.6:3.6,samples=100]
      (\x,{exp(-\x*\x/2)/sqrt(2*pi)});
    \draw[gold,very thick,dashed] (0,0)--(0,.44);
    \node[above,font=\scriptsize] at (0,.44) {$\mu$};
    \foreach \k/\l in {-3/\mu{-}3\sigma,-2/\mu{-}2\sigma,-1/\mu{-}\sigma,
        1/\mu{+}\sigma,2/\mu{+}2\sigma,3/\mu{+}3\sigma} {
        \draw (\k,0)--(\k,-.015);
        \node[below,font=\tiny,inner sep=1pt] at (\k,-.01) {$\l$};
      }
    \draw[navy,<->] (-1,-.1)--(1,-.1);
    \node[fill=white,inner sep=1pt,font=\tiny] at (0,-.1) {$68\%$};
    \draw[navy,<->] (-2,-.145)--(2,-.145);
    \node[fill=white,inner sep=1pt,font=\tiny] at (0,-.145) {$95\%$};
    \draw[navy,<->] (-3,-.19)--(3,-.19);
    \node[fill=white,inner sep=1pt,font=\tiny] at (0,-.19) {$99.7\%$};
  \end{tikzpicture}}
% The standard normal Phi(z) as the area left of z, and the symmetry
% Phi(-z) = 1 - Phi(z): the two gold tails have the same area.
\newcommand{\figphi}{%
  \begin{tikzpicture}[x=.3cm,y=3.2cm]
    \fill[teal!30] (-3.3,0) -- plot[domain=-3.3:1,samples=60]
      (\x,{exp(-\x*\x/2)/sqrt(2*pi)}) -- (1,0) -- cycle;
    \draw[->] (-3.5,0)--(3.6,0);
    \draw[very thick,navy] plot[domain=-3.3:3.3,samples=80]
      (\x,{exp(-\x*\x/2)/sqrt(2*pi)});
    \node[below,font=\scriptsize] at (1,0) {$z$};
    \node[below,font=\scriptsize] at (0,0) {$0$};
    \node[font=\tiny] at (-.6,.15) {$\Phi(z)$};
    \node[font=\tiny] at (0,.5) {area left of $z$};
    \begin{scope}[xshift=2.45cm]
      \fill[gold!50] (-3.3,0) -- plot[domain=-3.3:-1,samples=40]
        (\x,{exp(-\x*\x/2)/sqrt(2*pi)}) -- (-1,0) -- cycle;
      \fill[gold!50] (1,0) -- plot[domain=1:3.3,samples=40]
        (\x,{exp(-\x*\x/2)/sqrt(2*pi)}) -- (3.3,0) -- cycle;
      \draw[->] (-3.5,0)--(3.6,0);
      \draw[very thick,navy] plot[domain=-3.3:3.3,samples=80]
        (\x,{exp(-\x*\x/2)/sqrt(2*pi)});
      \node[below,font=\scriptsize] at (-1,0) {$-z$};
      \node[below,font=\scriptsize] at (1,0) {$z$};
      \node[font=\tiny] at (0,.5) {$\Phi(-z)=1-\Phi(z)$};
    \end{scope}
  \end{tikzpicture}}
% The unit square: the joint density of two variables on [0,1]^2 and a
% region A (below the diagonal) whose probability is a double integral.
\newcommand{\figjointsquare}{%
  \begin{tikzpicture}[scale=1.5]
    \fill[teal!12] (0,0) rectangle (1.4,1.4);
    \fill[teal!45] (0,0) -- (1.4,0) -- (1.4,1.4) -- cycle;
    \draw[->] (-.1,0) -- (1.7,0) node[right,font=\scriptsize] {$x$};
    \draw[->] (0,-.1) -- (0,1.7) node[above,font=\scriptsize] {$y$};
    \draw[thick,navy] (0,0) rectangle (1.4,1.4);
    \draw[navy,densely dashed] (0,0) -- (1.4,1.4);
    \node[below,font=\scriptsize] at (1.4,0) {$1$};
    \node[left,font=\scriptsize] at (0,1.4) {$1$};
    \node[font=\scriptsize] at (1.0,.35) {$A$};
    \node[font=\tiny,rotate=45] at (.55,.72) {$y=x$};
  \end{tikzpicture}}

% The exercises (one file each in exercises/) and \exerciseframe
% Exercises of Week 10, shared by beamer.tex, the study guide, the script and
% exercises.tex. Every exercise lives in a file of its own,
% exercises/NN-slug.tex, numbered 10.NN (see exercises.lua for the docstring).
% The body of an exercise file:
%
%   \begin{problem} ... \end{problem}      the statement, used everywhere;
%                                          inside it, \begin{context} ...
%                                          \end{context} is background that
%                                          the slides leave out
%   \begin{slideproblem} ... \end{slideproblem}
%                                          (optional) a shorter statement
%                                          for the slides
%   \begin{slides} ... \end{slides}        one frame of the deck: \exstatement
%                                          places the statement; repeat the
%                                          environment for more frames
%   \begin{solution} ... \end{solution}    the detailed solution (study
%                                          guide, exercise PDFs, website)
%
% Commands for the documents (the argument is the slug: 01-license-plates.tex
% has the slug license-plates):
%
%   \exref{slug}                 the number, e.g. 10.1, for "Exercise~\exref{...}"
%   \extitle{slug}               the title from the docstring
%   \exerciseframe[opts]{slug}   beamer: the frames of the exercise
%   \studyexercise{slug}         study guide: heading, statement, solution
%   \scriptexercise{slug}        script: the statement box, as on the slide
%   \allexercises                exercises.tex: every exercise, in order; on
%                                the website statement and solution with
%                                markers where the pipeline splits the HTML
%                                (\collectionexercise), in the PDF plain:
%                                each statement, followed by its solution if
%                                \ifsolutions (exercises-solutions.tex)

\def\exweek{10}

\directlua{exercises = dofile("exercises/exercises.lua")
  exercises.scan("exercises", "\exweek")}

\ExplSyntaxOn

\msg_new:nnn { exercises } { unknown } { Unknown~exercise~`#1'. }

\cs_new_protected:Npn \__ex_missing:n #1
  { \msg_error:nnn { exercises } { unknown } {#1} ?? }
\cs_new:Npn \exref #1
  {
    \cs_if_exist:cTF { exnum@#1 }
      { \use:c { exnum@#1 } } { \__ex_missing:n {#1} }
  }
\cs_new:Npn \extitle #1
  {
    \cs_if_exist:cTF { extitle@#1 }
      { \use:c { extitle@#1 } } { \__ex_missing:n {#1} }
  }

% ---------------------------------------------------------------------------
% Loading an exercise: its parts are stored, nothing is typeset. The slides
% typeset the stored parts; the other documents read the part they need from
% the file again (\__ex_input_part:nn), as stored tables do not survive lwarp.

\tl_new:N \g__ex_current_tl
\tl_new:N \g__ex_problem_tl
\tl_new:N \g__ex_slideproblem_tl
\tl_new:N \g__ex_solution_tl
\seq_new:N \g__ex_slides_seq
\box_new:N \l__ex_scratch_box

\cs_new_protected:Npn \__ex_load:n #1
  {
    \cs_if_exist:cF { exfile@#1 } { \__ex_missing:n {#1} }
    \tl_gset:Nn \g__ex_current_tl {#1}
    \tl_gclear:N \g__ex_problem_tl
    \tl_gclear:N \g__ex_slideproblem_tl
    \tl_gclear:N \g__ex_solution_tl
    \seq_gclear:N \g__ex_slides_seq
    \vbox_set:Nn \l__ex_scratch_box
      {
        \DeclareDocumentEnvironment { problem } { +b }
          { \tl_gset:Nn \g__ex_problem_tl {##1} } { }
        \DeclareDocumentEnvironment { slideproblem } { +b }
          { \tl_gset:Nn \g__ex_slideproblem_tl {##1} } { }
        \DeclareDocumentEnvironment { slides } { +b }
          { \seq_gput_right:Nn \g__ex_slides_seq {##1} } { }
        \DeclareDocumentEnvironment { solution } { +b }
          { \tl_gset:Nn \g__ex_solution_tl {##1} } { }
        \exp_args:Nv \input { exfile@#1 }
      }
  }

% Typesets one part (problem, slideproblem, solution) of an exercise straight
% from its file; the other parts are skipped.
\cs_new_protected:Npn \__ex_input_part:nn #1#2
  {
    \group_begin:
      \clist_map_inline:nn { problem, slideproblem, slides, solution }
        {
          \str_if_eq:nnTF {##1} {#2}
            { \DeclareDocumentEnvironment {##1} { } { } { } }
            { \DeclareDocumentEnvironment {##1} { +b } { } { } }
        }
      \exp_args:Nv \input { exfile@#1 }
    \group_end:
  }

% Exercise 10.1, or Exercise 10.2 · Finance
\cs_new:Npn \__ex_heading:
  {
    Exercise~\exref { \g__ex_current_tl }
    \cs_if_exist:cT { exfinance@ \g__ex_current_tl } { ~·~Finance }
  }
\prg_new_conditional:Npnn \__ex_if_finance: { T, F, TF }
  {
    \cs_if_exist:cTF { exfinance@ \g__ex_current_tl }
      { \prg_return_true: } { \prg_return_false: }
  }

% ---------------------------------------------------------------------------
% Beamer

% The finance context is printed in the study guide, the exercise PDFs and on
% the website; the
% slides and the script (which shows what the slide shows) leave it out.
\bool_new:N \l__ex_context_bool
\bool_set_true:N \l__ex_context_bool
\DeclareDocumentEnvironment { context } { }
  {
    \bool_if:NTF \l__ex_context_bool
      { \par \textit { Finance~context. } ~ }
      { \__ex_gobble_context:w }
  }
  { \bool_if:NT \l__ex_context_bool { \par } }
\cs_new_protected:Npn \__ex_gobble_context:w #1 \end #2
  { \str_if_eq:nnTF {#2} { context } { \end { context } } { \__ex_gobble_context:w } }

\IfClassLoadedTF { beamer }
  {
    \bool_set_false:N \l__ex_context_bool

    \cs_new:Npn \__ex_block: { \__ex_if_finance:TF { exampleblock } { block } }
    \NewDocumentCommand \exstatement { }
      {
        \begin { \__ex_block: } { \__ex_heading: }
          \tl_if_empty:NTF \g__ex_slideproblem_tl
            { \g__ex_problem_tl } { \g__ex_slideproblem_tl }
        \end { \__ex_block: }
      }

    \tl_new:N \l__ex_frames_tl
    \NewDocumentCommand \exerciseframe { O{t} m }
      {
        \__ex_load:n {#2}
        \tl_clear:N \l__ex_frames_tl
        \seq_map_indexed_inline:Nn \g__ex_slides_seq
          {
            \tl_put_right:Nn \l__ex_frames_tl { \begin{frame}[#1] }
            \tl_put_right:Ne \l__ex_frames_tl
              {
                {
                  \exp_not:N \extitle {#2}
                  \int_compare:nNnT {##1} > { 1 } { \c_space_tl (cont.) }
                }
              }
            \tl_put_right:Nn \l__ex_frames_tl { ##2 \end{frame} }
          }
        \tl_use:N \l__ex_frames_tl
      }
  }
  {
    \cs_new:Npn \__ex_box: { \__ex_if_finance:TF { fin } { task } }
    \cs_new_protected:Npn \__ex_statement_box:n #1
      {
        \begin { \__ex_box: } { \__ex_heading: }
          \__ex_input_part:nn { \g__ex_current_tl } {#1}
        \end { \__ex_box: }
      }

    \NewDocumentCommand \studyexercise { m }
      {
        \__ex_load:n {#1}
        \subsection* { Exercise~\exref{#1}~---~\extitle{#1} }
        \addcontentsline { toc } { subsection }
          { Exercise~\exref{#1}~---~\extitle{#1} }
        \__ex_statement_box:n { problem }
        \soltitle \__ex_input_part:nn {#1} { solution }
      }

    \NewDocumentCommand \scriptexercise { m }
      {
        \__ex_load:n {#1}
        \group_begin:
          \bool_set_false:N \l__ex_context_bool
          \tl_if_empty:NTF \g__ex_slideproblem_tl
            { \__ex_statement_box:n { problem } }
            { \__ex_statement_box:n { slideproblem } }
        \group_end:
      }

    % The website (lwarp): statement and solution of each exercise, between
    % markers where the pipeline splits the HTML.
    \NewDocumentCommand \collectionexercise { m }
      {
        \__ex_load:n {#1}
        \par [[exercise:#1:problem]] \par
        \__ex_input_part:nn {#1} { problem }
        \par [[exercise:#1:solution]] \par
        \__ex_input_part:nn {#1} { solution }
        \par [[exercise:#1:end]] \par
      }
    % The PDF: plain, each statement, followed by its solution in the
    % version with solutions (\ifsolutions). The statement is set off like a
    % quote in Markdown, by a bar on the left (quote=<color>, also used by
    % exercises.tex) in the color of its box in the study guide: navy, or teal
    % for a finance exercise.
    \tcbset
      {
        quote/.style =
          {
            enhanced, breakable, frame~hidden, interior~hidden,
            sharp~corners, boxrule = 0pt, borderline~west = {2pt}{0pt}{#1},
            left = 8pt, right = 0pt, top = 1pt, bottom = 1pt,
            before~skip = 6pt, after~skip = 6pt,
            % the document's paragraph spacing (parskip=half)
            before~upper = { \setlength { \parskip } { .5\baselineskip } }
          }
      }
    \NewDocumentCommand \plainexercise { m }
      {
        \__ex_load:n {#1}
        \subsection* { Exercise~\exref{#1}~---~\extitle{#1} }
        \addcontentsline { toc } { subsection }
          { Exercise~\exref{#1}~---~\extitle{#1} }
        \__ex_if_finance:TF
          { \begin{tcolorbox}[quote=teal] }
          { \begin{tcolorbox}[quote=navy] }
          \__ex_input_part:nn {#1} { problem }
        \end{tcolorbox}
        \legacy_if:nT { solutions }
          {
            \soltitle
            \__ex_input_part:nn {#1} { solution }
          }
      }
    \NewDocumentCommand \allexercises { }
      {
        \IfPackageLoadedTF { lwarp }
          { \directlua { exercises.each("collectionexercise") } }
          { \directlua { exercises.each("plainexercise") } }
      }
  }

\ExplSyntaxOff


% PDF properties and the date on the title slide
% The properties of every PDF of Week 10: author (with email address), title,
% subject, keywords and the date the document last changed, in the PDF
% (hyperref) and on its title page (\date), and the header and footer of the
% printed documents (below). Each document loads this file after hyperref and
% then calls \documentinfo{<what it is>}, e.g. \documentinfo{Slides}.
%
% The date: the website build (scripts/build-content.mjs) puts
% \def\lastmodified{YYYY-MM-DD} in front of the document, the date of the last
% commit that changed its sources. Compiled locally, it is today.

\newcommand{\courseinfo}{Mathematics for Finance}
\newcommand{\weeklabel}{Week 10}
\newcommand{\weekinfo}{\weeklabel: Continuous random variables, joint distributions, independence}
\newcommand{\authorname}{Tibor SÁNDOR}
\providecommand{\myemail}{tiborsandor2001@gmail.com}

\makeatletter
\ifdefined\lastmodified\else
	\edef\lastmodified{\the\year-\two@digits{\the\month}-\two@digits{\the\day}}
\fi
\def\di@split#1-#2-#3\@nil{%
	\def\di@year{#1}\def\di@month{#2}\def\di@day{#3}}
\expandafter\di@split\lastmodified\@nil
% November 11, 2026
\edef\lastmodifiedtext{%
	\ifcase\numexpr\di@month\relax\or January\or February\or March\or April\or
		May\or June\or July\or August\or September\or October\or November\or
		December\fi\space\number\di@day, \di@year}
% D:20261111000000, the date format of PDF metadata
\edef\di@pdfdate{D:\di@year\di@month\di@day000000}

\newcommand{\documentinfo}[1]{%
	\hypersetup{
		pdftitle={\courseinfo, \weekinfo\space -- #1},
		pdfauthor={\authorname\space<\myemail>},
		pdfsubject={\courseinfo, \weekinfo},
		pdfkeywords={probability, random variables, mathematics for finance,
				Corvinus University of Budapest},
		pdfcreationdate={\di@pdfdate},
		pdfmoddate={\di@pdfdate},
		pdflang={en-US},
	}%
	\date{Last updated: \lastmodifiedtext}%
	\gdef\documentkind{#1}%
}

% Header and footer of the printed documents (KOMA-Script, scrlayer-scrpage;
% the slides have beamer's footline, the website neither): the course, the
% week and the document on the left of the header (the week's topic is on the
% title page), the university on the right; in the footer the author and the
% page as <current>/<total>. The title page gets only the footer.
\newcommand{\documentkind}{}
\IfClassLoadedTF{beamer}{}{\IfPackageLoadedTF{lwarp}{}{%
	\usepackage{lastpage}%
	\usepackage[headsepline=1.5pt,footsepline=1.5pt,plainfootsepline]
		{scrlayer-scrpage}%
	\clearpairofpagestyles
	\setkomafont{pageheadfoot}{\small\sffamily\color{navy}}%
	\setkomafont{pagenumber}{}%
	\setkomafont{headsepline}{\color{navy!35}}%
	\setkomafont{footsepline}{\color{navy!35}}%
	\ihead{\courseinfo\ \textperiodcentered\ \weeklabel\ \textperiodcentered\ 
		\documentkind}%
	\ohead{Corvinus University of Budapest}%
	\ifoot*{\authorname\ \textperiodcentered\ \myemail}%
	\ofoot*{\pagemark/\pageref*{LastPage}}%
	\pagestyle{scrheadings}%
	% Extra stretch for breaking lines with long inline formulas.
	\setlength{\emergencystretch}{1.5em}%
}}
\makeatother

\documentinfo{Slides}

\title[Mathematics for Finance]{Mathematics for Finance}
\author{Tibor SÁNDOR}
\subtitle{Week 10: Continuous random variables, named continuous
  distributions, joint distributions, independence}
\institute{Corvinus University of Budapest\\[2pt]
  \href{mailto:\myemail}{\texttt{\myemail}}}

\begin{document}

\begin{frame}
  \titlepage
\end{frame}

\begin{frame}{Plan for today}
  \begin{enumerate}
    \item \textbf{Session 1 (90 min):} continuous random variables and
          their density; the distribution function; expected value and
          variance; the named continuous distributions: uniform,
          exponential, normal.
    \item \textbf{Session 2 (90 min):} the distribution of a function of a
          random variable; joint distributions of two random variables,
          discrete and continuous; the joint distribution function;
          independent random variables.
  \end{enumerate}
  \vspace{.3cm}
  \textbf{Why does a finance student need this?} Waiting times, losses and
  returns can take any value in an interval, and the risk of a portfolio
  depends on how its parts move \emph{together}.
  \vspace{.2cm}

  Each topic: the key ideas, then an exercise. More exercises, all with
  full solutions, are in the exercises PDF for self-study.
\end{frame}

%======================================================================
\section{Session 1: Continuous random variables, named distributions}
%======================================================================

\begin{frame}{Session 1: Continuous random variables and named
    distributions}
  \textbf{Last week:} discrete random variables. A pmf $p(x)=P(X=x)$, the
  distribution function $F(x)=P(X\le x)$, $\E[X]=\sum_x x\,p(x)$,
  $\Var(X)$, and five named distributions: Bernoulli, binomial, Poisson,
  geometric, hypergeometric.
  \vspace{3mm}

  \textbf{Today:} the same questions for a random variable that can take
  \emph{any} value in an interval.
  \begin{itemize}
    \item the density $f$ and the distribution function $F$;
    \item expected value, variance, standard deviation: sums become
          integrals;
    \item three named continuous distributions: \textbf{uniform},
          \textbf{exponential}, \textbf{normal}.
  \end{itemize}
\end{frame}

\begin{frame}{From discrete to continuous}
  \small
  \begin{columns}[c]
    \begin{column}{.55\textwidth}
      \textbf{Discrete:} $X$ takes the values $x_1,x_2,\dots$; each value
      has a probability $p(x_i)=P(X=x_i)$, and they add up to 1.
      \vspace{2mm}

      \textbf{Continuous:} the time until the next order, the loss on a
      loan, a return. These can take \emph{any} value in an interval.
      \begin{itemize}
        \item There are too many values to give each a positive
              probability: $P(X=a)=0$ for every single number $a$.
        \item We ask for the probability of an \emph{interval}:
              $P(a\le X\le b)$.
        \item It is the \emph{area} under a curve, the \textbf{density}
              $f$.
      \end{itemize}
    \end{column}
    \begin{column}{.43\textwidth}
      \centering
      \figdensity
    \end{column}
  \end{columns}
\end{frame}

\begin{frame}{Continuous random variables and their density}
  \small
  \begin{block}{Definition}
    $X$ is \textbf{continuous} if there is a function $f$, its
    \textbf{density}, such that for every interval
    \[ P(a\le X\le b)=\int_a^b f(x)\,dx . \]
  \end{block}
  \begin{columns}[T]
    \begin{column}{.55\textwidth}
      \begin{block}{A density has two properties}
        (1) $f(x)\ge 0$ for every $x$;\\[2pt]
        (2) $\displaystyle\int_{-\infty}^{\infty} f(x)\,dx=1$: the total
        area is 1.
      \end{block}
    \end{column}
    \begin{column}{.43\textwidth}
      \begin{itemize}
        \item $P(X=a)=\int_a^a f=0$, so
              $P(a\le X\le b)=P(a<X<b)$.
        \item $f(x)$ is \emph{not} a probability: it can be bigger
              than 1. Only areas are probabilities.
      \end{itemize}
    \end{column}
  \end{columns}
\end{frame}

\exerciseframe{settlement-time}

\begin{frame}{The distribution function}
  \small
  \begin{columns}[T]
    \begin{column}{.53\textwidth}
      \begin{block}{Distribution function (any random variable)}
        \[ F(x)=P(X\le x). \]
        For a continuous $X$: \
        $F(x)=\displaystyle\int_{-\infty}^{x} f(t)\,dt$, \ and \
        $f(x)=F'(x)$.
      \end{block}
      \begin{block}{Probabilities from $F$}
        $P(a<X\le b)=F(b)-F(a)$, \quad $P(X>a)=1-F(a)$.
      \end{block}
      $F$ grows from 0 to 1; for a continuous $X$ it has no jumps (a
      discrete $X$ has a step function).
    \end{column}
    \begin{column}{.45\textwidth}
      \centering
      \figcdf\\[1mm]
      \footnotesize the density of the first picture: the area between
      $a$ and $b$ is the rise of $F$
    \end{column}
  \end{columns}
\end{frame}

\exerciseframe{loss-given-default}

\begin{frame}{Expected value and variance}
  \small
  Sums become integrals, $p(x)$ becomes $f(x)$:
  \begin{block}{Expected value}
    \[ \E[X]=\int_{-\infty}^{\infty} x\,f(x)\,dx,\qquad
      \E[g(X)]=\int_{-\infty}^{\infty} g(x)\,f(x)\,dx,\qquad
      \E[aX+b]=a\,\E[X]+b. \]
  \end{block}
  \begin{block}{Variance and standard deviation}
    With $\mu=\E[X]$: \
    $\Var(X)=\E\big[(X-\mu)^2\big]=\E[X^2]-\mu^2$, \
    $\SD(X)=\sqrt{\Var(X)}$,\\[2pt]
    $\Var(aX+b)=a^2\Var(X)$, \ $\SD(aX+b)=|a|\,\SD(X)$.
  \end{block}
  \textbf{Example.} $f(x)=2x$ on $[0,1]$ (0 outside):
  \[ \E[X]=\int_0^1 x\cdot 2x\,dx=\frac23,\qquad
    \E[X^2]=\int_0^1 x^2\cdot 2x\,dx=\frac12,\qquad
    \Var(X)=\frac12-\frac49=\frac1{18}. \]
\end{frame}

\exerciseframe{recovery-rate}

\begin{frame}{The uniform distribution}
  \small
  \begin{block}{Uniform distribution on $[\alpha,\beta]$}
    \[ f(x)=\begin{cases}\dfrac{1}{\beta-\alpha}, & \alpha\le x\le\beta,\\
        0, & \text{otherwise,}\end{cases}\qquad
      F(x)=\frac{x-\alpha}{\beta-\alpha}\ \ (\alpha\le x\le\beta). \]
    $\E[X]=\dfrac{\alpha+\beta}{2}$, \qquad
    $\Var(X)=\dfrac{(\beta-\alpha)^2}{12}$.
  \end{block}
  \begin{columns}[c]
    \begin{column}{.5\textwidth}
      \begin{itemize}
        \item A number ``picked at random'' from $[\alpha,\beta]$: every
              interval of the same length is equally likely.
        \item $P(c\le X\le d)=\dfrac{d-c}{\beta-\alpha}$: length over total
              length.
        \item $F=0$ left of $\alpha$, $F=1$ right of $\beta$.
      \end{itemize}
    \end{column}
    \begin{column}{.48\textwidth}
      \centering
      \figuniform
    \end{column}
  \end{columns}
\end{frame}

\exerciseframe{waiting-for-the-bus}

\begin{frame}{The exponential distribution}
  \small
  \begin{block}{Exponential distribution with rate $\lambda>0$}
    \[ f(x)=\begin{cases}\lambda e^{-\lambda x}, & x\ge0,\\
        0, & x<0,\end{cases}\qquad
      F(x)=1-e^{-\lambda x},\quad P(X>x)=e^{-\lambda x}\ \ (x\ge0). \]
    $\E[X]=\dfrac1\lambda$, \qquad $\Var(X)=\dfrac1{\lambda^2}$.
  \end{block}
  \begin{columns}[c]
    \begin{column}{.5\textwidth}
      \begin{itemize}
        \item Waiting times: the next order, the next claim, a default.
              $\lambda$ = events per unit of time.
        \item \textbf{Memoryless:} for $s,t\ge0$
              \[ P(X>s+t\mid X>s)=P(X>t). \]
              Having waited $s$ does not shorten the rest of the wait.
      \end{itemize}
    \end{column}
    \begin{column}{.48\textwidth}
      \centering
      \figexponential
    \end{column}
  \end{columns}
\end{frame}

\exerciseframe{order-arrivals}

\begin{frame}{The normal distribution}
  \small
  \begin{columns}[T]
    \begin{column}{.53\textwidth}
      \begin{block}{Normal distribution $N(\mu,\sigma^2)$}
        \[ f(x)=\frac{1}{\sigma\sqrt{2\pi}}\,
          e^{-(x-\mu)^2/(2\sigma^2)},\qquad x\in\mathbb{R}. \]
        $\E[X]=\mu$, \ $\Var(X)=\sigma^2$: the bell is centered at $\mu$,
        $\sigma$ is its width.
      \end{block}
      \begin{block}{Standard normal $Z\sim N(0,1)$}
        $\Phi(z)=P(Z\le z)$, from the table; \ $\Phi(-z)=1-\Phi(z)$.
      \end{block}
      \textbf{Standardizing:} if $X\sim N(\mu,\sigma^2)$, then
      $Z=\dfrac{X-\mu}{\sigma}\sim N(0,1)$, so
      \[ P(X\le x)=\Phi\Big(\frac{x-\mu}{\sigma}\Big). \]
    \end{column}
    \begin{column}{.45\textwidth}
      \centering
      \fignormal\\[1mm]
      \footnotesize the 68--95--99.7 rule: $X$ is within $\sigma$,
      $2\sigma$, $3\sigma$ of $\mu$ with probability about $68\%$, $95\%$,
      $99.7\%$
    \end{column}
  \end{columns}
\end{frame}

\begin{frame}{The standard normal table}
  \footnotesize
  \begin{columns}[c]
    \begin{column}{.62\textwidth}
      \centering
      $\Phi(z)=P(Z\le z)$ for $Z\sim N(0,1)$\\[1.5mm]
      \renewcommand{\arraystretch}{1.08}
      \begin{tabular}{@{}cc@{\hspace{5mm}}cc@{\hspace{5mm}}cc@{\hspace{5mm}}cc@{}}
        \toprule
        $z$ & $\Phi(z)$ & $z$ & $\Phi(z)$ & $z$ & $\Phi(z)$ & $z$ & $\Phi(z)$\\
        \midrule
        0.0 & 0.5000 & 0.8 & 0.7881 & 1.6 & 0.9452 & 2.4 & 0.9918\\
        0.1 & 0.5398 & 0.9 & 0.8159 & 1.7 & 0.9554 & 2.5 & 0.9938\\
        0.2 & 0.5793 & 1.0 & 0.8413 & 1.8 & 0.9641 & 2.6 & 0.9953\\
        0.3 & 0.6179 & 1.1 & 0.8643 & 1.9 & 0.9713 & 2.7 & 0.9965\\
        0.4 & 0.6554 & 1.2 & 0.8849 & 2.0 & 0.9772 & 2.8 & 0.9974\\
        0.5 & 0.6915 & 1.3 & 0.9032 & 2.1 & 0.9821 & 2.9 & 0.9981\\
        0.6 & 0.7257 & 1.4 & 0.9192 & 2.2 & 0.9861 & 3.0 & 0.9987\\
        0.7 & 0.7580 & 1.5 & 0.9332 & 2.3 & 0.9893 &     &       \\
        \bottomrule
      \end{tabular}\\[2.5mm]
      \textbf{Quantiles} $z_p$ with $\Phi(z_p)=p$:\\[1mm]
      \begin{tabular}{@{}c*{5}{c}@{}}
        \toprule
        $p$   & 0.90  & 0.95  & 0.975 & 0.99  & 0.995\\
        $z_p$ & 1.282 & 1.645 & 1.960 & 2.326 & 2.576\\
        \bottomrule
      \end{tabular}
    \end{column}
    \begin{column}{.36\textwidth}
      \centering
      \figphi\\[2mm]
      \raggedright
      \begin{itemize}
        \item Negative $z$: \ $\Phi(-z)=1-\Phi(z)$.
        \item $P(Z>z)=1-\Phi(z)$.
        \item $P(a<Z\le b)=\Phi(b)-\Phi(a)$.
        \item $P(|Z|\le z_{0.975})=0.95$.
      \end{itemize}
    \end{column}
  \end{columns}
\end{frame}

\exerciseframe{annual-return}

\begin{frame}{Session 1: summary}
  \small
  \begin{itemize}
    \item Density: $P(a\le X\le b)=\int_a^b f(x)\,dx$; \ $f\ge0$, total
          area 1, $P(X=a)=0$.
    \item $F(x)=P(X\le x)=\int_{-\infty}^x f(t)\,dt$; \ $f=F'$; \
          $P(a<X\le b)=F(b)-F(a)$.
    \item $\E[X]=\int x f(x)\,dx$, \ $\E[g(X)]=\int g(x)f(x)\,dx$, \
          $\Var(X)=\E[X^2]-(\E[X])^2$, \ $\SD(X)=\sqrt{\Var(X)}$.
  \end{itemize}
  \vspace{1mm}
  \centering
  \renewcommand{\arraystretch}{1.25}
  \begin{tabular}{@{}llll@{}}
    \toprule
    Distribution & density $f(x)$ & $\E[X]$ & $\Var(X)$\\\midrule
    uniform on $[\alpha,\beta]$ & $\frac{1}{\beta-\alpha}$ on
    $[\alpha,\beta]$ & $\frac{\alpha+\beta}{2}$ &
    $\frac{(\beta-\alpha)^2}{12}$\\
    exponential $(\lambda)$ & $\lambda e^{-\lambda x}$, $x\ge0$; \
    $P(X>x)=e^{-\lambda x}$ & $\frac1\lambda$ & $\frac1{\lambda^2}$\\
    normal $N(\mu,\sigma^2)$ & $\frac{1}{\sigma\sqrt{2\pi}}
    e^{-(x-\mu)^2/(2\sigma^2)}$; \
    $P(X\le x)=\Phi\big(\frac{x-\mu}{\sigma}\big)$ & $\mu$ &
    $\sigma^2$\\\bottomrule
  \end{tabular}
  \vspace{1.5mm}

  \raggedright
  \textbf{Self-study:} Exercises~\exref{device-lifetime},
  \exref{from-cdf-to-density}, \exref{power-density},
  \exref{default-time}, \exref{standard-normal-table}.
\end{frame}

%======================================================================
\section{Session 2: Functions of random variables, joint distributions}
%======================================================================

\begin{frame}{Session 2: Functions of random variables, joint
    distributions, independence}
  \small
  \textbf{Recap of Session 1}
  \begin{itemize}
    \item Density: $P(a\le X\le b)=\int_a^b f(x)\,dx$; \
          $F(x)=P(X\le x)$, \ $f=F'$, \ $P(a<X\le b)=F(b)-F(a)$.
    \item $\E[X]=\int x f(x)\,dx$, \ $\Var(X)=\E[X^2]-(\E[X])^2$.
    \item Uniform: length over total length. Exponential:
          $P(X>x)=e^{-\lambda x}$. Normal: $P(X\le x)=
          \Phi\big(\frac{x-\mu}{\sigma}\big)$.
  \end{itemize}
  \vspace{2mm}
  \textbf{Now}
  \begin{itemize}
    \item the distribution of $Y=g(X)$, e.g.\ a stock price $S=S_0e^X$;
    \item two random variables at the same time: how do two stocks, two
          loans, two claims behave \emph{together}?
  \end{itemize}
\end{frame}

\begin{frame}{The distribution of a function of a random variable}
  \small
  We know the distribution of $X$; what is the distribution of
  $Y=g(X)$, e.g.\ $Y=X^2$ or $Y=S_0e^X$?
  \begin{block}{Method: first the distribution function}
    (1) $F_Y(y)=P(Y\le y)=P(g(X)\le y)$; \
    (2) rewrite $g(X)\le y$ as an event about $X$ and use the distribution
    of $X$; \
    (3) if a density is needed: $f_Y(y)=F_Y'(y)$.
  \end{block}
  \textbf{Example.} $X$ uniform on $[0,1]$, $Y=2X+3$. For $3\le y\le 5$:
  \[ F_Y(y)=P(2X+3\le y)=P\Big(X\le \frac{y-3}{2}\Big)=\frac{y-3}{2}, \qquad
    f_Y(y)=\frac12 , \]
  so $Y$ is uniform on $[3,5]$.
  \begin{block}{A linear function of a normal is normal}
    If $X\sim N(\mu,\sigma^2)$ and $a\ne0$, then
    $aX+b\sim N(a\mu+b,\,a^2\sigma^2)$. \ In particular
    $\dfrac{X-\mu}{\sigma}\sim N(0,1)$.
  \end{block}
\end{frame}

\exerciseframe{square-of-uniform}

\exerciseframe{lognormal-price}

\begin{frame}{Joint distribution: the discrete case}
  \small
  \begin{columns}[T]
    \begin{column}{.55\textwidth}
      \begin{block}{Joint pmf}
        $p(x,y)=P(X=x,\,Y=y)$: a table; all entries add up to 1.
      \end{block}
      \begin{block}{Marginal pmfs}
        $p_X(x)=\sum_y p(x,y)$ \ (row sums),\\[2pt]
        $p_Y(y)=\sum_x p(x,y)$ \ (column sums).
      \end{block}
      \begin{block}{Joint distribution function (any $X$, $Y$)}
        $F(a,b)=P(X\le a,\,Y\le b)=\displaystyle\sum_{x\le a}\sum_{y\le b}
        p(x,y)$
      \end{block}
    \end{column}
    \begin{column}{.43\textwidth}
      \centering
      \renewcommand{\arraystretch}{1.25}
      \begin{tabular}{c|cc|c}
        $X\backslash Y$ & $0$ & $1$ & $p_X$\\\hline
        $0$ & $0.3$ & $0.2$ & \alert{$0.5$}\\
        $1$ & $0.1$ & $0.4$ & \alert{$0.5$}\\\hline
        $p_Y$ & \alert{$0.4$} & \alert{$0.6$} & $1$
      \end{tabular}\\[2mm]
      \footnotesize the sums are written in the \emph{margins}: hence
      ``marginal''\\[1mm]
      here $F(0,1)=0.3+0.2=0.5$, \ $F(0,0)=0.3$
    \end{column}
  \end{columns}
\end{frame}

\exerciseframe{two-stocks}

\begin{frame}{Joint distribution: the continuous case}
  \small
  \begin{columns}[T]
    \begin{column}{.6\textwidth}
      \begin{block}{Joint density}
        $f(x,y)\ge0$, total $\iint f=1$, and for a region $A$
        \[ P\big((X,Y)\in A\big)=\iint_A f(x,y)\,dx\,dy . \]
      \end{block}
      \begin{block}{Marginal densities: integrate the other variable out}
        $f_X(x)=\int_{-\infty}^{\infty} f(x,y)\,dy$, \quad
        $f_Y(y)=\int_{-\infty}^{\infty} f(x,y)\,dx$
      \end{block}
      \begin{block}{Joint distribution function}
        $F(a,b)=P(X\le a,\,Y\le b)=\int_{-\infty}^{a}
        \!\int_{-\infty}^{b} f(x,y)\,dy\,dx$
      \end{block}
    \end{column}
    \begin{column}{.38\textwidth}
      \centering
      \figjointsquare\\[1mm]
      \footnotesize the probability of a region $A$ (here $X>Y$) is the
      volume under $f$ above it; do the inner integral first
    \end{column}
  \end{columns}
\end{frame}

\exerciseframe{joint-density-square}

\begin{frame}{Independent random variables}
  \small
  \begin{block}{Definition}
    $X$ and $Y$ are \textbf{independent} if for all sets $A$, $B$
    \[ P(X\in A,\ Y\in B)=P(X\in A)\,P(Y\in B), \]
    equivalently $F(a,b)=F_X(a)\,F_Y(b)$ for all $a$, $b$.
  \end{block}
  \begin{columns}[T]
    \begin{column}{.49\textwidth}
      \begin{block}{Discrete}
        $p(x,y)=p_X(x)\,p_Y(y)$ for \emph{every} cell of the table.
        One cell that fails: not independent.
      \end{block}
    \end{column}
    \begin{column}{.49\textwidth}
      \begin{block}{Continuous}
        $f(x,y)=f_X(x)\,f_Y(y)$ for all $x$, $y$.\\
        On a rectangle: $f$ splits into
        (function of $x$)$\times$(function of $y$).
      \end{block}
    \end{column}
  \end{columns}
  \vspace{1mm}
  Knowing the value of $X$ does not change the distribution of $Y$.
  Probabilities of ``and'' multiply.
\end{frame}

\exerciseframe{two-stocks-independence}

\exerciseframe{two-default-times}

\exerciseframe{meeting-time}

\begin{frame}{Session 2: summary}
  \small
  \begin{itemize}
    \item $Y=g(X)$: first $F_Y(y)=P(g(X)\le y)$, then $f_Y=F_Y'$. \
          $X\sim N(\mu,\sigma^2)\Rightarrow aX+b\sim N(a\mu+b,a^2\sigma^2)$.
    \item Discrete: joint pmf $p(x,y)$ in a table; \ $p_X(x)=\sum_y
          p(x,y)$, \ $p_Y(y)=\sum_x p(x,y)$.
    \item Continuous: $P((X,Y)\in A)=\iint_A f(x,y)\,dx\,dy$; \
          $f_X(x)=\int f(x,y)\,dy$, \ $f_Y(y)=\int f(x,y)\,dx$.
    \item Joint distribution function $F(a,b)=P(X\le a,\,Y\le b)$.
    \item Independent: $p(x,y)=p_X(x)p_Y(y)$, or $f(x,y)=f_X(x)f_Y(y)$,
          for \emph{all} $x$, $y$; then $P(X\in A,Y\in B)=P(X\in A)P(Y\in
          B)$.
    \item The marginals alone do not tell how $X$ and $Y$ move together.
  \end{itemize}
  \vspace{1mm}
  \textbf{Self-study:} Exercises~\exref{log-of-uniform},
  \exref{three-coins}, \exref{exponential-pair},
  \exref{which-are-independent}.
\end{frame}

\end{document}
